% ------------------------------------------------------------
% §6  The research program, cut into undergraduate-sized pieces
% ------------------------------------------------------------
\section{Undergraduate short projects}\label{sec:projects}

The preceding sections flagged thirteen concrete projects at the point where
each becomes natural; here we collect them and add the program-level view.
The unifying activity is \emph{measurement}: locating scaled-down members
of the \SHA{} design family in coordinate systems (Flajolet's above all)
where ``looks random'' becomes a number with error bars. None of the
projects requires cryptographic background beyond this paper; all require
programming and a taste for experiment; each produces figures and data
that are a contribution regardless of what they show. That last property
is worth underlining for research with undergraduates in short windows:
these are experiments that \emph{cannot fail}, only surprise.

\begin{center}
\small
\begin{tabular}{@{}p{0.30\linewidth}p{0.22\linewidth}p{0.12\linewidth}p{0.24\linewidth}@{}}
  \toprule
  project (defined in) & prerequisites & window & deliverable \\
  \midrule
  Single ops under the $2$-adic lens (\cref{sec:statespace}) &
    elementary number theory & 3--4 wks &
    verified ergodicity criteria, orbit visualizations \\[2pt]
  The linear algebra of diffusion (\cref{sec:rivest}) &
    linear algebra & 3--4 wks &
    invertibility \& group structure of the $\Sigma/\sigma$ maps \\[2pt]
  Exact inverses of $\Sigma_0, \Sigma_1$ (\cref{sec:rivest}) &
    polynomials over $\Ftwo$ & 2--3 wks &
    closed-form inverse maps, verified two ways \\[2pt]
  Functional-graph atlas of truncated \SHA{} (\cref{sec:telescope}) &
    probability, programming & 6--8 wks &
    atlas of measured vs.\ predicted statistics \\[2pt]
  Collision hunting, van Oorschot--Wiener style (\cref{sec:telescope}) &
    probability, programming & 4--6 wks &
    timing distributions vs.\ theory; parallel implementation \\[2pt]
  The dynamics of iteration (\cref{sec:iterated}) &
    probability, programming & 6--8 wks &
    image-shrinkage law, fixed points, Hellman trade-off \\[2pt]
  Cutoff for toy compression walks (\cref{sec:prngs}) &
    probability, linear algebra & 4--6 wks &
    spectral gaps and cutoff profiles vs.\ rounds \\[2pt]
  Cayley hashes by hand (\cref{sec:provable}) &
    group theory, linear algebra & 3--4 wks &
    spectral gap vs.\ equidistribution; collision gallery \\[2pt]
  Differential trails in miniature (\cref{sec:siege}) &
    probability, programming & 4--6 wks &
    exact toll tables; best-trail probability vs.\ rounds \\[2pt]
  The \textsc{sat} cliff (\cref{sec:npsearch}) &
    programming, logic & 4--6 wks &
    solver-time cliff chart; three frontiers on one axis \\[2pt]
  A neural distinguisher at toy scale (\cref{sec:mlcrypt}) &
    programming, machine learning & 6--8 wks &
    learned-advantage vs.\ rounds; interpretation of the model \\[2pt]
  Where do the tests start failing? (\cref{sec:gap}) &
    statistics, programming & 4--6 wks &
    round-by-round battery failure chart \\[2pt]
  Avalanche cartography (\cref{sec:gap}) &
    linear algebra, programming & 4--6 wks &
    influence-matrix heatmaps by round \\
  \bottomrule
\end{tabular}
\end{center}

\medskip

Three remarks on the program as a whole.

\paragraph{Model organisms.} The full-scale function is beyond experiment
(\cref{sec:telescope}) and beyond proof (\cref{sec:gap}); the program
therefore studies the \emph{family} obtained by shrinking each design
parameter---word size $32 \to 8$ or $16$, state words $8 \to 4$, rounds
$64 \to r$---the way biology studies fruit flies. The two-geometry anatomy
of \cref{sec:statespace} says which shrinkings are honest: one must
preserve the ARX alternation, or the object degenerates into something a
single lens solves. A systematic scaled family, with its statistics
measured, would be a reusable community resource; no standard one exists.

\paragraph{The comparative method.} Every measurement gains meaning from
controls on both sides: genuinely random maps (the Flajolet silhouette)
on one side, and visibly structured maps (bijections, $\Ftwo$-linear
maps, one-round toys) on the other. The interesting science is in the
\emph{trajectory}---how, as rounds accumulate, a \SHA{}-family map walks
from the structured region onto the random silhouette, in which
coordinates it arrives first, and whether the walk is gradual or a phase
transition. SHA-1 and \SHA{} can be given the identical treatment; given
\cref{sec:falling}, any measurable divergence between their trajectories
would be a finding of genuine interest.

\paragraph{Odds of surprise.} Realistically: most measurements will
confirm the random silhouette, and their value is the atlas itself plus
trained students. But the history of \cref{sec:josephus} cuts both
ways---structure, when it exists in innocent-looking discrete processes,
has repeatedly been found by exactly this kind of patient small-case
empiricism, and essentially nobody has pointed the analytic-combinatorics
instrument at this family at model-organism scale. The expected value of
looking is small-positive; the tail is long.

\medskip

\noindent\emph{On tooling.} Everything above runs comfortably in Python
or Julia on a laptop for $n \le 30$; the collision-hunting and atlas
projects scale gracefully to whatever cluster time is available. We
maintain reference implementations and are happy to share scaffolding so
that student time goes to science rather than plumbing.
